% sections/classification.tex

\section{Spill Classification}
\label{sec:classification}

Every spill is classified into one of three tiers \emph{before} any cleanup begins. Classification
determines who responds, what PPE is required, and whether the building must be evacuated. When in
doubt, classify upward (treat as the more serious tier) and call the EHS emergency line.

\begin{tierbox}{spillteal}{TIER 1 -- Minor Spill}
\textbf{Definition:} Small quantity (typically under 1 liter or 1 kg), a single known low-hazard
chemical, no fire/reactivity/inhalation risk, and within the training and PPE of the person present.

\textit{Examples:} spilled dilute buffer, small volume of ethanol, a broken vial of table salt
solution.

\textbf{Who responds:} The lab occupant who discovered the spill, or any trained lab member
present, using the local spill kit.
\end{tierbox}

\vspace{0.35em}

\begin{tierbox}{spillamber}{TIER 2 -- Moderate Spill}
\textbf{Definition:} Larger quantity (roughly 1--4 liters or 1--4 kg), a flammable or corrosive
liquid, a spill that has spread beyond an immediate bench area, or one that exceeds the
responder's training or the kit's capacity.

\textit{Examples:} a fallen bottle of concentrated hydrochloric acid, several liters of a
flammable solvent, a mercury thermometer break with visible beading.

\textbf{Who responds:} Lab Supervisor and EHS Office; area is isolated and ventilated while
awaiting EHS arrival.
\end{tierbox}

\vspace{0.35em}

\begin{tierbox}{spillred}{TIER 3 -- Major / Emergency Spill}
\textbf{Definition:} Any spill involving a highly toxic, highly reactive, or unknown chemical;
a spill accompanied by fire, smoke, or a strong/irritating vapor cloud; a spill that has caused
injury or exposure; or any quantity that cannot be safely contained by trained lab personnel.

\textit{Examples:} a broken container of concentrated nitric acid reacting with organic waste,
a cryogenic liquid spill causing an oxygen-deficiency alarm, an unlabeled unknown liquid spreading
across a hallway.

\textbf{Who responds:} Building evacuation, Campus Emergency Line (ext.\ 5911), Fire Department
Hazmat Team. No lab personnel attempt cleanup.
\end{tierbox}
